\section{Block recoding}\label{sec:recursive-addressing}

The only recoding needed later is the standard block presentation that
turns a finite union of cylinders into a one-step Markov hole.

\begin{lemma}[Finite observation under block recoding]
\label{thm:finite-depth-collapse}
Let \(Y\subseteq A^{\NNzero}\) be a one-sided subshift and let
\(\Psi_\ell:Y\to Y^{(\ell)}\) be the usual \(\ell\)-block code,
\[
  \Psi_\ell(y)_n:=y_n\cdots y_{n+\ell-1}.
\]
If \(\mathcal F_m\) denotes the observation \(\sigma\)-algebra generated
by the first \(m\) original symbols and \(\mathcal F_m^{(\ell)}\) the
one generated by the first \(m\) block symbols, then
\[
  \Psi_\ell^{-1}(\mathcal F_m^{(\ell)})
  \subseteq
  \mathcal F_{m+\ell-1}.
\]
Consequently, for every event \(\mathsf P\) that is measurable with
respect to the block coding,
\[
  \varepsilon_{m+\ell-1}(\mathsf P;\mu)
  \leq
  \varepsilon_m(\mathsf P;\Psi_{\ell*}\mu).
\]
\end{lemma}

\begin{proof}
The first \(m\) block symbols are determined by the first \(m+\ell-1\)
original symbols, so every \(\mathcal F_m^{(\ell)}\)-measurable set
pulls back to an \(\mathcal F_{m+\ell-1}\)-measurable set.  The Bayes
error can only decrease when one conditions on a finer
\(\sigma\)-algebra.
\end{proof}
